\chapter*{Preface to the first edition}
\addcontentsline{toc}{chapter}{Preface to the first edition}
\markboth{Preface to the first edition}{Preface to the first edition}
\readingsource{Lawrence C. Evans: original preface}%
  {English printed pp.~xiii--xv; PDF pages~14--16.
   Chinese reference: PDF pages~11--12.\newline
   原著序言全文；中文边注为整理者补充。保留作者的写作原则、致谢和署名。}

% BEGIN ORIGINAL PREFACE
I present in this book a wide-ranging survey of many important topics in the
theory of partial differential equations (PDE), with particular emphasis on
various modern approaches.%
\PDEVocab{survey; exposition}{survey 在此指综述；exposition 指系统讲解。全书既讲显式解公式，也讲能量方法、变分法与守恒律。}
I have made a huge number of editorial decisions
about what to keep and what to toss out, and can only claim that this selection
seems to me about right. I of course include the usual formulas for solutions
of the usual linear PDE, but also devote large amounts of exposition to energy
methods within Sobolev spaces, to the calculus of variations, to conservation
laws, etc.

My general working principles in the writing have been these:

\noindent\textbf{a. PDE theory is (mostly) not restricted to two independent variables.}
\PDEIdea{作者提醒：偏微分方程不限于两个自变量，也研究高维问题。他质疑的是把二维分类当作普适框架，并非否定具体方程的分类。}
Many texts describe PDE as if functions of the two variables $(x,y)$ or $(x,t)$
were all that matter. This emphasis seems to me misleading, as modern
discoveries concerning many types of equations, both linear and nonlinear,
have allowed for the rigorous treatment of these in any number of dimensions.
I also find it unsatisfactory to ``classify'' partial differential equations:
this is possible in two variables, but creates the false impression that there
is some kind of general and useful classification scheme available in general.

\noindent\textbf{b. Many interesting equations are nonlinear.}
\PDEVocab{accordingly; analogues}{accordingly：因此、相应地。nonlinear analogues 指线性理论在非线性问题中的相应思想或结果；能否推广仍须逐项证明。}
My view is that overall we know too much about linear PDE and too little about
nonlinear PDE. I have accordingly introduced nonlinear concepts early in the
text and have tried hard to emphasize everywhere nonlinear analogues of the
linear theory.

\clearpage
\noindent\textbf{c. Understanding generalized solutions is fundamental.}
\PDEIdea{“弱化解的定义”与“保持唯一性”必须一起考虑。弱解、熵解、黏性解有各自的定义和适用情形，不能简单混用。}
Many of the partial differential equations we study, especially nonlinear
first-order equations, do not in general possess smooth solutions. It is
therefore essential to devise some kind of proper notion of generalized or
weak solution. This is an important but subtle undertaking, and much of the
hardest material in this book concerns the uniqueness of appropriately
defined weak solutions.

\noindent\textbf{d. PDE theory is not a branch of functional analysis.}
\PDEChinese{作者强调不能只依靠抽象算子框架；具体的微积分计算、测度论和估计同样关键。这并不排斥使用泛函分析。}
Whereas certain classes of equations can profitably be viewed as generating
abstract operators between Banach spaces, the insistence on an overly abstract
viewpoint, and consequent ignoring of deep calculus and measure theoretic
estimates, is ultimately limiting.

\noindent\textbf{e. Notation is a nightmare.}
\PDEVocab{at variance with}{意为“与……不一致”。本书为保持全书统一，部分符号可能不同于某个专门领域的惯例；阅读时先核对附录 A。}
I have really tried to introduce consistent notation, which works for all the
important classes of equations studied. This attempt is sometimes at variance
with notational conventions within a given subarea.

\noindent\textbf{f. Good theory is (almost) as useful as exact formulas.}
\PDEChinese{三部分的主线：先求显式公式，再研究线性方程的一般理论，最后研究若干重要非线性方程。存在性、唯一性等结论本身也是“求解”的成果。}
I incorporate this principle into the overall organization of the text, which
is subdivided into three parts, roughly mimicking the historical development
of PDE theory itself. Part~I concerns the search for explicit formulas for
solutions, and Part~II the abandoning of this quest in favor of general theory
asserting the existence and other properties of solutions for linear equations.
Part~III is the mostly modern endeavor of fashioning general theory for
important classes of nonlinear PDE.

Let me also explicitly comment here that I intend the development within each
section to be rigorous and complete (exceptions being the frankly heuristic
treatment of asymptotics in \S4.5 and an occasional reference to a research
paper).%
\PDEVocab{heuristic / sharp}{heuristic：启发式的；arcane：深奥的；sharp：条件、常数等更精确或接近最优。作者优先解释思想；初读可暂缓部分难证明，不能把启发式推导当作完整证明。}
This means that even locally within each chapter the topics do not
necessarily progress logically from ``easy'' to ``hard'' concepts. There are
many difficult proofs and computations early on, but as compensation many
easier ideas later. The student should certainly omit on first reading some
of the more arcane proofs.

I wish next to emphasize that this is a textbook, and not a reference book.
I have tried everywhere to present the essential ideas in the clearest possible
settings, and therefore have almost never established sharp versions of any
of the theorems. Research articles and advanced monographs, many of them
listed in the Bibliography, provide such precision and generality. My goal
has rather been to explain, as best I can, the many fundamental ideas of the
subject within fairly simple contexts.

I have greatly profited from the comments and thoughtful suggestions of many
of my colleagues, friends and students, in particular: S.~Antman, J.~Bang,
X.~Chen, A.~Chorin, M.~Christ, J.~Cima, P.~Colella, J.~Cooper, M.~Crandall,
B.~Driver, M.~Feldman, M.~Fitzpatrick, R.~Gariepy, J.~Goldstein, D.~Gomes,
O.~Hald, W.~Han, W.~Hrusa, T.~Ilmanen, I.~Ishii, I.~Israel, R.~Jerrard,
C.~Jones, B.~Kawohl, S.~Koike, J.~Lewis, T.-P.~Liu, H.~Lopes, J.~McLaughlin,
K.~Miller, J.~Morford, J.~Neu, M.~Portilheiro, J.~Ralston, F.~Rezakhanlou,
W.~Schlag, D.~Serre, P.~Souganidis, J.~Strain, W.~Strauss, M.~Struwe,
R.~Temam, B.~Tvedt, J.-L.~Vazquez, M.~Weinstein, P.~Wolfe, and Y.~Zheng.

I especially thank Tai-Ping Liu for many years ago writing out for me the
first draft of what is now Chapter~11.

I am extremely grateful for the suggestions and lists of mistakes from earlier
drafts of this book sent to me by many readers, and I encourage others to send
me their comments, at \nolinkurl{evans@math.berkeley.edu}. I have come to realize
that I must be more than slightly mad to try to write a book of this length
and complexity, but I am not yet crazy enough to think that I have made no
mistakes. I will therefore maintain a listing of errors which come to light,
and will make this accessible through the \texttt{math.berkeley.edu} homepage.

Faye Yeager at UC Berkeley has done a really magnificent job typing and
updating these notes, and Jaya Nagendra heroically typed an earlier version at
the University of Maryland. My deepest thanks to both.

I have been supported by the NSF during much of the writing, most recently
under grant DMS-9424342.

\begin{flushright}
LCE\\
August, 1997\\
Berkeley
\end{flushright}
% END ORIGINAL PREFACE
